\chapter{Objetivos: Geral e Específicos}
\label{chap:objetivos}


A pesquisa proposta busca abordar as fraudes identificadas em sistemas de gerenciamento, com foco na emissão de GTA e nas transações relacionadas à movimentação de bovinos. Para isso, propõe-se o desenvolvimento e adaptação de um sistema neuro-fuzzy evolutivo, capaz de tratar dados complexos, dinâmicos e incertos, ajustando-se continuamente a novos padrões e estratégias de fraude. Essa abordagem visa detectar e prevenir fraudes de forma mais eficiente, promovendo maior integridade, transparência e confiabilidade ao sistema.

\section{Objetivo Geral}

Propor, adaptar e implementar novos algoritmos baseados em sistemas inteligentes evolutivos com métodos de pertinência fuzzy para detecção e prevenção de fraudes em sistemas de gerenciamento agropecuário, com foco na adaptabilidade, autonomia e interpretabilidade dos modelos propostos.

\section{Objetivos Específicos}
\begin{enumerate}
  \item \textbf{Realizar revisão abrangente da literatura sobre:} 
  \begin{itemize}
    \item Sistemas inteligentes evolutivos e adaptativos
    \item Sistemas neuro-fuzzy evolutivos e métodos de pertinência
    \item Técnicas de detecção de anomalias e fraudes em sistemas de gestão
    \item Algoritmos de aprendizado incremental e online
    \item Metodologias de adaptação estrutural e paramétrica em tempo real
  \end{itemize}
  
  \item \textbf{Propor e adaptar novos algoritmos evolutivos:} 
  \begin{itemize}
    \item Desenvolver algoritmos com capacidade de adaptação estrutural automática
    \item Implementar mecanismos de aprendizado incremental para dados de fluxo contínuo
    \item Integrar métodos de pertinência fuzzy para modelagem de incertezas e ambiguidades
    \item Garantir interpretabilidade e explicabilidade dos modelos propostos
  \end{itemize}
  
  \item \textbf{Adaptar técnicas ao contexto agropecuário:} 
  \begin{itemize}
    \item Personalizar algoritmos para características específicas de dados agropecuários
    \item Desenvolver métodos de tratamento de dados esparsos e ruidosos
    \item Implementar técnicas de detecção de drift conceitual em ambientes dinâmicos
  \end{itemize}
  
  \item \textbf{Validar algoritmos com dados reais:} 
  \begin{itemize}
    \item Testar modelos com datasets históricos de sistemas agropecuários
    \item Alcançar métricas de desempenho superiores aos métodos estado da arte
    \item Demonstrar capacidade de adaptação em cenários de mudança conceitual
  \end{itemize}
  
  \item \textbf{Avaliar e comparar desempenho:} 
  \begin{itemize}
    \item Comparar algoritmos propostos com métodos de referência usando métricas padronizadas
    \item Realizar análises estatísticas de significância dos resultados
    \item Avaliar eficiência computacional e escalabilidade dos métodos
  \end{itemize}
  
  \item \textbf{Publicar contribuições científicas:} 
  \begin{itemize}
    \item Produzir artigos em periódicos de alto impacto na área
    \item Apresentar resultados em conferências nacionais e internacionais
    \item Disponibilizar código e datasets para reprodutibilidade científica
  \end{itemize}
\end{enumerate}

\section{Resultados Esperados}

\begin{enumerate}
  \item Modelo inteligente para detecção de fraudes em sistemas de gerenciamento do setor agropecuário.
  \item Base de dados organizada e anonimizada para estudos futuros, com uma estrutura crescente de 8.298.490 registros de GTAs validados e estruturados.
  \item Relatórios sobre padrões fraudulentos e propostas de mitigação baseados em análise de dados reais.
  \item Publicações científicas em detecção de fraudes no setor agropecuário com validação experimental robusta.
  \item Propostas de melhorias que aumentem a eficiência dos sistemas de gerenciamento do setor agropecuário.
  \item Impacto positivo no setor, com redução de fraudes, fortalecimento da segurança alimentar e impactos na arrecadação.
  \item Algoritmos validados com base de dados crescente de 8.298.490 registros estruturados demonstrando aplicabilidade prática.
\end{enumerate}